\chapter{Sign sequences: from experiments to a research boundary}\label{ch:signs}
\goal{Measure how closely a sequence of plus and minus signs matches shifted copies of itself, and translate the question into one about matrices whose rows are orthogonal. Prove two conditions that the size of such a matrix must satisfy, compare them with a complete computer search over short sequences, and see why a search can never replace a proof that covers every length. Finally, read a recent research claim about these matrices and say exactly what it adds to what you have proved.}

\section{A sequence with small echoes}
A sign sequence of positive length $n$ is a list $a=(a_0,\ldots,a_{n-1})$ with each entry equal to $1$ or $-1$. Its aperiodic autocorrelation at a shift $s$ with $1\le s<n$ is
\[
A_s(a)=\sum_{j=0}^{n-1-s}a_j a_{j+s}.
\]
Slide the sequence by $s$ positions, multiply overlapping entries, and add the products. ``Aperiodic'' means entries beyond the endpoints are not wrapped around. At shift zero, the analogous sum is $n$, because every sign has square one.

For $a=(1,1,1,-1)$,
\begin{align*}
A_1&=(1)(1)+(1)(1)+(1)(-1)=1,\\
A_2&=(1)(1)+(1)(-1)=0,\\
A_3&=(1)(-1)=-1.
\end{align*}
A Barker sequence is a sign sequence with $|A_s|\le1$ at every nonzero aperiodic shift. Thus this example qualifies. The definition makes sense at length one as well: there are no nonzero shifts to check, so the condition holds vacuously.

Every $A_s$ is a sum of $n-s$ signs and therefore has the same parity as $n-s$. To see this, if $r$ summands are negative, the sum is $(n-s-r)-r=n-s-2r$. Hence a Barker sequence must have $A_s=0$ when $n-s$ is even, and $A_s=1$ or $-1$ when $n-s$ is odd. The definition has immediately produced an arithmetic restriction.

\subsection*{Decode the summation limits before computing}
At shift $s$, a term $a_j a_{j+s}$ is available only while both indices belong to $0,1,\ldots,n-1$. Starting at $j=0$, the last allowed $j$ satisfies $j+s=n-1$, or $j=n-1-s$. There are therefore $n-s$ terms, including both endpoints of that index range.

For a length-five sequence and shift two, the sum is
\[
 A_2=a_0a_2+a_1a_3+a_2a_4.
\]
There is no term $a_3a_5$ in an aperiodic correlation, because $a_5$ is outside the list. At the largest allowed shift $s=n-1$, only $a_0a_{n-1}$ remains. Checking these endpoint cases is an effective way to detect an incorrect summation limit.

The subscript on $A_s$ records the shift, whereas subscripts on $a_j$ record positions. The uppercase quantity is a sum built from the lowercase entries. Keeping these roles distinct helps when several indices appear in one calculation.

\section{Periodic correlation is a different quantity}
For periodic correlation, wrap the indices around modulo $n$:
\[
C_s(a)=\sum_{j=0}^{n-1}a_j a_{(j+s)\bmod n},\qquad0\le s<n.
\]
There are always $n$ summands. For $a=(1,1,1,-1)$, direct multiplication gives $C_0=4$ and $C_1=C_2=C_3=0$.

For $1\le s<n$, the split occurs precisely at $j=n-s$. When $0\le j\le n-1-s$, the index $j+s$ is still below $n$, and the contribution is $A_s$. When $n-s\le j\le n-1$, the index wraps once and becomes $j+s-n$. Thus
\[
 C_s=\sum_{j=0}^{n-1-s}a_j a_{j+s}
     +\sum_{j=n-s}^{n-1}a_j a_{j+s-n}.
\]
In the second sum, set $r=j+s-n$. At its first term $j=n-s$, we have $r=0$; at its last term $j=n-1$, we have $r=s-1$. Solving the same equation for $j$ gives $j=r+n-s$. Substitution therefore yields
\begin{align*}
 \sum_{j=n-s}^{n-1}a_j a_{j+s-n}
 &=\sum_{r=0}^{s-1}a_{r+n-s}a_r\\
 &=\sum_{r=0}^{s-1}a_r a_{r+n-s}\\
 &=A_{n-s}.
\end{align*}
The middle equality commutes the two real factors. The final equality uses the definition of $A_{n-s}$: its last index is $n-1-(n-s)=s-1$. This is a change of indexing, not an approximation.

For a concrete check, take $n=5,s=2$. The periodic sum has five terms:
\[
 C_2=\underbrace{a_0a_2+a_1a_3+a_2a_4}_{A_2}
       +\underbrace{a_3a_0+a_4a_1}_{A_3}.
\]
The last two become $a_0a_3+a_1a_4$ on reversing the factors. The example displays exactly the general split. Therefore
\[
C_s=A_s+A_{n-s}.
\tag{14.1}\label{eq:periodaperiod}
\]
The two definitions are related but not interchangeable. For $a=(1,1,-1)$, we have $A_1=0$, $A_2=-1$, so it is Barker; yet $C_1=-1$, not zero. A claim about one kind of correlation cannot be transferred to the other without a justified bridge.

\section{Circulant matrices package the same products}\label{ch14:sec:circulant}
The periodic correlations of a sequence can be read off from a matrix. Write the sequence $a$ as the first row. Obtain each further row from the one above it by moving every entry one place to the right, with the last entry wrapping around to the front. After $n$ rows, one more step would bring back the first row, so we stop there and have an $n$-by-$n$ matrix. To match the positions in the sequence, we number its rows and columns from $0$ to $n-1$, rather than from $1$ to $n$ as in Chapter~\ref{ch:linear}.

\par\noindent\begin{minipage}{\linewidth}
\begin{definition}[Circulant matrix]\label{ch14:def:circulant}
Let $a$ be a sign sequence of length $n$. The \emph{circulant matrix} of $a$ is the $n$-by-$n$ matrix $H$ with entries
\[
H_{ij}=a_{(j-i)\bmod n},\qquad0\le i,j\le n-1.
\]
A matrix of this form is called \emph{circulant}.
\end{definition}
\end{minipage}\par

Row $i$ is the sequence rotated $i$ places to the right. For instance, in row $1$ the entry in column $0$ is $a_{(0-1)\bmod n}=a_{n-1}$: the last entry of the sequence has wrapped around to the front. Row $0$ is the sequence itself, because $H_{0j}=a_{j\bmod n}=a_j$. So every circulant matrix is the circulant matrix of its own first row. For our length-four example $a=(1,1,1,-1)$,
\[
H=\begin{pmatrix}
1&1&1&-1\\
-1&1&1&1\\
1&-1&1&1\\
1&1&-1&1
\end{pmatrix}.
\]
Each row is the row above it with its last entry moved to the front.

As $j$ runs through $0,1,\ldots,n-1$, the index $(j-i)\bmod n$ also runs through $0,1,\ldots,n-1$, each value exactly once, only starting at a different place. For instance, with $n=4$ and $i=1$, the columns $j=0,1,2,3$ give the indices $3,0,1,2$. In the language of Section~\ref{ch03:sec:injsurj}, the map $j\mapsto(j-i)\bmod n$ is a bijection from $\{0,1,\ldots,n-1\}$ to itself, and its inverse is $k\mapsto(k+i)\bmod n$. So each row of $H$ contains every entry of $a$ exactly once, in a rotated order. In particular, every row has squared length $a_0^2+a_1^2+\cdots+a_{n-1}^2=n$.

\subsection*{Dot products of rows are periodic correlations}
Now compare two rows of $H$, row $i$ and row $\ell$. Their dot product (Section~\ref{ch10:sec:inner}) is
\[
 \sum_{j=0}^{n-1}a_{(j-i)\bmod n}\,a_{(j-\ell)\bmod n}.
\]
To simplify it, rename the summation index: put $k=(j-i)\bmod n$. We have just seen that $k$ runs through $0,1,\ldots,n-1$, each value once, as $j$ does. So we may sum over $k$ instead of $j$; the terms are the same, only listed in a different order. The second factor must also be written in terms of $k$. Since $k\equiv j-i\pmod n$,
\[
 j-\ell=(j-i)+(i-\ell)\equiv k+(i-\ell)\pmod n.
\]
Congruent integers have the same remainder on division by $n$ (Section~\ref{ch06:sec:residues}), so $(j-\ell)\bmod n=(k+i-\ell)\bmod n$. Now let $s=(i-\ell)\bmod n$. Then $s\equiv i-\ell$, so $(k+i-\ell)\bmod n=(k+s)\bmod n$ for the same reason, and the dot product becomes
\[
 \sum_{k=0}^{n-1}a_k\,a_{(k+s)\bmod n}=C_s.
\]

\par\noindent\begin{minipage}{\linewidth}
\begin{proposition}\label{ch14:prop:rows}
Let $a$ be a sign sequence of length $n$, and let $H$ be its circulant matrix.
\begin{enumerate}
\item[(a)] For all rows $i$ and $\ell$, the dot product of row $i$ with row $\ell$ equals $C_{(i-\ell)\bmod n}(a)$.
\item[(b)] Any two different rows of $H$ are orthogonal exactly when $C_s(a)=0$ for every shift $s$ with $1\le s\le n-1$.
\end{enumerate}
\end{proposition}
\end{minipage}\par

\begin{proof}
Part (a) is the calculation above. For part (b) we need to know which shifts occur in part (a). If $i\ne\ell$, then $i-\ell$ lies between $-(n-1)$ and $n-1$ and is not zero, so it is not a multiple of $n$, and the shift $(i-\ell)\bmod n$ is not zero. Conversely, every shift $s$ with $1\le s\le n-1$ does occur: rows $i=s$ and $\ell=0$ give $(s-0)\bmod n=s$.

Now suppose that $C_s=0$ for every $s$ with $1\le s\le n-1$. By the first observation, the dot product of two different rows is one of these correlations, so it is zero. Conversely, suppose that any two different rows are orthogonal. By the second observation, each $C_s$ with $1\le s\le n-1$ is the dot product of the two different rows $s$ and $0$, so it is zero.
\end{proof}

Both observations are needed. The first gives the ``if'' direction of part~(b), and the second gives the ``only if'' direction. Neither one alone would show that the two conditions are equivalent.

\subsection*{Read the matrix identity entry by entry}
Section~\ref{ch10:sec:product} called a square matrix with entries $\pm1$ a \emph{Hadamard matrix} when it satisfies $HH^{\mathsf T}=nI$, where $n$ is the number of rows. This number is called the \emph{order} of the matrix. The identity $HH^{\mathsf T}=nI$ contains $n^2$ equations between numbers, one for each position $(i,\ell)$. As we saw in that section, the entry of $HH^{\mathsf T}$ in position $(i,\ell)$ is the dot product of row $i$ with row $\ell$. The entry of $nI$ in that position is $n$ if $i=\ell$ and zero otherwise. So the equations on the diagonal say that every row has squared length $n$, and the equations off the diagonal say that different rows are orthogonal. For a matrix with entries $\pm1$, the diagonal equations hold automatically, because each row consists of $n$ entries with square one. So such a matrix is a Hadamard matrix exactly when any two different rows are orthogonal.

For the matrix $H$ above, rows $0$ and $1$ have dot product
\[
(1)(-1)+(1)(1)+(1)(1)+(-1)(1)=-1+1+1-1=0.
\]
By Proposition~\ref{ch14:prop:rows}(a) this is $C_{(0-1)\bmod4}=C_3$, which we computed to be zero in Section~\ref{ch14:sec:periodic}. Row $0$ dotted with itself gives $1+1+1+1=4$. The proposition takes care of all the remaining pairs at once. A matrix identity is useful because it collects many related equations into one line, but to understand it you have to know what a single one of its entries says.

\subsection*{Circulant Hadamard matrices}
A \emph{circulant Hadamard matrix} is a Hadamard matrix that is also circulant. By Proposition~\ref{ch14:prop:rows}(b), the circulant matrix of $a$ is a Hadamard matrix exactly when $C_s(a)=0$ for every shift $s$ with $1\le s\le n-1$. Since every circulant matrix is the circulant matrix of its own first row, we arrive at the following statement:
\begin{quote}
A circulant Hadamard matrix of order $n$ exists exactly when some sign sequence of length $n$ has $C_s=0$ for every shift $s$ with $1\le s\le n-1$.
\end{quote}
Our example $(1,1,1,-1)$ has $C_1=C_2=C_3=0$, so the matrix $H$ displayed above is a circulant Hadamard matrix of order four.

Habit~6 (Section~\ref{ch06:sec:order}) asks what a change of description preserves. Here nothing is lost in the translation: a question about the correlations of a sequence and a question about the rows of a matrix have turned out to be the same question in two languages. Each language makes different tools available. The sequence language suggests sums and shifts, while the matrix language suggests dot products and orthogonality. The next two sections use both.

\section{The sum of all correlations reveals a square}
Let $r=\sum_{j=0}^{n-1}a_j$ be the row sum. Add every periodic correlation, including shift zero. By definition, each $C_s$ is itself a sum over $j$, so the total is a double sum. A finite sum may be added in any order, so we may add over $s$ first for each fixed $j$, and then pull out the factor $a_j$, which does not depend on $s$:
\begin{align*}
\sum_{s=0}^{n-1}C_s
&=\sum_{s=0}^{n-1}\sum_{j=0}^{n-1}a_j a_{(j+s)\bmod n}
 =\sum_{j=0}^{n-1}\sum_{s=0}^{n-1}a_j a_{(j+s)\bmod n}\\
&=\sum_{j=0}^{n-1}a_j\left(\sum_{s=0}^{n-1}a_{(j+s)\bmod n}\right)\\
&=\sum_{j=0}^{n-1}a_j\left(\sum_{k=0}^{n-1}a_k\right)\\
&=r^2.
\end{align*}
In the second line, for each fixed $j$, the indices $(j+s)\bmod n$ run through every position exactly once. This is the specific symmetry that justifies the reindexing.

One way to see the whole count is to use three abstract entries $a_0,a_1,a_2$. The three periodic correlations list these products:
\begin{align*}
 C_0&=a_0a_0+a_1a_1+a_2a_2,\\
 C_1&=a_0a_1+a_1a_2+a_2a_0,\\
 C_2&=a_0a_2+a_1a_0+a_2a_1.
\end{align*}
Together they contain each ordered product $a_j a_k$ once. Expanding $(a_0+a_1+a_2)^2$ gives those same nine ordered products, so the totals agree. For general $n$, each fixed first index $j$ meets every second index exactly once as the shift varies. That is why the total is $(\sum_j a_j)^2$, rather than, for example, $\sum_j a_j^2$.

This identity itself does not require signs; the same finite rearrangement holds for real entries. The restriction to signs is used afterward to identify $C_0=n$ and to make the row sum an integer. Separating those roles tells us which part of the argument might generalize.

If all nonzero $C_s$ vanish, the left side is just $C_0=n$. Hence
\[
n=r^2.
\tag{14.2}\label{eq:squareorder}
\]
The integer $r$ is a sum of signs, so the order $n$ must be a square. This is a proved necessary condition, valid for every positive length, not a pattern extrapolated from an experiment.

\section{Orthogonality also forces divisibility by four}\label{ch14:sec:four}
The square condition of Section~\ref{ch14:sec:square} relied on the circulant pattern. The next condition holds for every Hadamard matrix, circulant or not, as long as its order is greater than $2$. Small orders show why that restriction is needed. The $1$-by-$1$ matrix $(1)$ is a Hadamard matrix of order $1$, and the matrix
\[
\begin{pmatrix}1&1\\1&-1\end{pmatrix}
\]
of Section~\ref{ch10:sec:product} is a Hadamard matrix of order $2$, but neither $1$ nor $2$ is divisible by four.

\par\noindent\begin{minipage}{\linewidth}
\begin{lemma}\label{thm:hadamardorder}
If a Hadamard matrix has order $n>2$, then $4$ divides $n$.
\end{lemma}
\end{minipage}\par

\begin{proof}
The idea is to replace the matrix by a simpler Hadamard matrix of the same order, and then to look closely at three of its rows. Two kinds of change keep a matrix Hadamard.
\begin{itemize}
\item \emph{Multiplying a column by $-1$.} This changes the sign of one entry in every row. In the dot product of two rows, the term that comes from that column is a product of two entries that have both changed sign, so it is unchanged, because $(-x)(-y)=xy$. The other terms are not touched at all. So every dot product of two rows stays the same, and the entries are still $\pm1$.
\item \emph{Rearranging the columns.} This applies the same rearrangement to every row. The dot product of two rows then consists of the same terms as before, added in a different order, so it does not change.
\end{itemize}
Neither change alters the order $n$, so it is enough to prove that $4$ divides $n$ for the new matrix. Even if the original matrix was circulant, the new one usually is not. That does not matter, because the lemma is about all Hadamard matrices, and the proof uses only two facts: the entries are $\pm1$, and different rows are orthogonal.

\emph{Step 1: make the first row consist of ones.} For each column whose entry in the first row is $-1$, multiply that column by $-1$. Afterward the first row is $(1,1,\ldots,1)$.

\emph{Step 2: the second row has as many plus signs as minus signs.} The second row is orthogonal to the first. Since the first row consists of ones, their dot product is simply the sum of the entries of the second row, so this sum is zero. If the second row has $p$ entries equal to $1$, then it has $n-p$ entries equal to $-1$, and its sum is $p-(n-p)=2p-n$. So $2p=n$. Hence $n$ is even, and the second row has $n/2$ plus signs and $n/2$ minus signs. Put $h=n/2$.

\emph{Step 3: sort the second row.} Rearrange the columns so that the $h$ columns in which the second row has a $1$ come first, and the $h$ columns in which it has a $-1$ come last. The first row still consists of ones, and the second row is now
\[
(\underbrace{1,\ldots,1}_{h},\underbrace{-1,\ldots,-1}_{h}).
\]

\emph{Step 4: look at a third row.} Since $n>2$, the matrix has a third row. Let $x$ be the sum of its first $h$ entries and $y$ the sum of its last $h$ entries. Its dot product with the first row is $x+y$. Its dot product with the second row is $x-y$, because the second row multiplies the first $h$ entries by $1$ and the last $h$ entries by $-1$. The third row is orthogonal to both, so
\[
x+y=0\qquad\text{and}\qquad x-y=0.
\]
Adding the two equations gives $2x=0$, so $x=0$. (Subtracting them gives $y=0$ as well, although we will not need it.)

\emph{Step 5: count again.} So the first $h$ entries of the third row are signs whose sum is zero. The count of Step~2 applies to them: if $q$ of these $h$ entries equal $1$, their sum is $q-(h-q)=2q-h$, and this is zero, so $h=2q$. Therefore $n=2h=4q$, and $4$ divides $n$.
\end{proof}

To see the proof in action, apply it to the circulant Hadamard matrix $H$ of Section~\ref{ch14:sec:circulant}. Its first row is $(1,1,1,-1)$, so Step~1 multiplies the last column by $-1$:
\[
\begin{pmatrix}
1&1&1&-1\\
-1&1&1&1\\
1&-1&1&1\\
1&1&-1&1
\end{pmatrix}
\quad\longrightarrow\quad
\begin{pmatrix}
1&1&1&1\\
-1&1&1&-1\\
1&-1&1&-1\\
1&1&-1&-1
\end{pmatrix}.
\]
The second row is now $(-1,1,1,-1)$, with its plus signs in columns $1$ and $2$ (counting from $0$). Step~3 moves these two columns to the front, so the new column order is $1,2,0,3$:
\[
\begin{pmatrix}
1&1&1&1\\
1&1&-1&-1\\
-1&1&1&-1\\
1&-1&1&-1
\end{pmatrix}.
\]
The first row consists of ones, and the second row is $(1,1,-1,-1)$. The third row is $(-1,1,1,-1)$. Its first half $(-1,1)$ and its second half $(1,-1)$ both add up to zero, as Step~4 predicts. Each half contains one plus sign, so $h=2\cdot1$ is even and $n=4\cdot1$ is divisible by four. You can check that the rows of the final matrix are still orthogonal. Notice also that it is no longer circulant: its second row is not its first row moved one place to the right.

Now we can combine the two necessary conditions.

\par\noindent\begin{minipage}{\linewidth}
\begin{proposition}\label{thm:circulantnecessary}
If a circulant Hadamard matrix has order $n$, then $n=1$ or $n=4u^2$ for some positive integer $u$.
\end{proposition}
\end{minipage}\par

\begin{proof}
Let $a$ be the first row of the matrix. The matrix is the circulant matrix of $a$, so by Proposition~\ref{ch14:prop:rows}(b), $C_s(a)=0$ for every shift $s$ with $1\le s\le n-1$. Equation~\eqref{eq:squareorder} therefore gives $n=r^2$, where $r$ is the row sum, an integer.

If $n=1$, there is nothing more to prove. Otherwise $n\ge2$. Moreover $n\ne2$, because $2$ is not the square of an integer: $1^2=1<2<4=2^2$. Hence $n>2$, and Lemma~\ref{thm:hadamardorder} shows that $4$ divides $n=r^2$. In particular $r^2$ is even, and so $r$ is even. Indeed, if $r$ were odd, say $r=2k+1$ for an integer $k$, then $r^2=4k^2+4k+1$ would be odd. Write $|r|=2u$, where $u$ is a nonnegative integer. Since $n=r^2$ is positive, $r\ne0$, and so $u\ge1$. Then $n=r^2=|r|^2=4u^2$.
\end{proof}

The proposition allows the orders $1$ and $4u^2$ for $u=1,2,3,\ldots$, that is,
\[
1,\ 4,\ 16,\ 36,\ 64,\ 100,\ \ldots.
\]
For the first two there are examples: the $1$-by-$1$ matrix $(1)$, and the matrix $H$ of order four. For the orders $16,36,64,\ldots$ the proposition says nothing either way. It neither builds a matrix of such an order nor rules one out. This is typical of a necessary condition: it removes most orders from consideration and tells us precisely which ones remain open. That list of remaining orders is the boundary between what our elementary arguments can decide and what needs new ideas.

\section{An exact finite experiment}\label{ch14:sec:experiment}
A sign sequence of length $n$ is fixed by $n$ independent choices, one sign for each position, so there are $2^n$ sign sequences of length $n$. For small $n$ a computer can examine every one of them. The table below shows the result of such an exhaustive search for $1\le n\le13$. Only integers are involved, so the arithmetic is exact, with no rounding. Appendix~\ref{app:practical} contains a short program that produces the table directly from the definitions. For each length the search counts two kinds of sequences:
\begin{itemize}
\item the Barker sequences;
\item the sequences with $C_s=0$ for every shift $s$ with $1\le s\le n-1$, which by Section~\ref{ch14:sec:circulant} are exactly the first rows of circulant Hadamard matrices.
\end{itemize}
Every list is counted on its own. For instance, a sequence and its negative are counted as two sequences. So are a sequence and its reversal (the same entries in the opposite order), and a sequence and its rotations, whenever these lists are different.

\begin{center}
\begin{tabular}{@{}r rr@{}}\toprule
Length $n$ & Barker sequences & Sequences with $C_s=0$ for $1\le s\le n-1$\\\midrule
1&2&2\\2&4&0\\3&4&0\\4&8&8\\5&4&0\\6&0&0\\7&4&0\\8&0&0\\9&0&0\\10&0&0\\11&4&0\\12&0&0\\13&4&0\\\bottomrule
\end{tabular}
\end{center}

The table is evidence about these thirteen lengths and no others. It cannot rule out a qualifying sequence of some larger length that was never tested. Nor can the search simply be pushed much further, because each additional position doubles the number of sequences. Length $60$ alone has $2^{60}$ sign sequences, which is more than $10^{18}$. A computer that examined a billion sequences per second would need more than $36$ years for that single length. More computing power is therefore not a way to settle the question for every length. Only a proof can do that.

Every pattern in the table agrees with what we have proved, and each one has a reason.
\begin{itemize}
\item The last column is zero at every length except $1$ and $4$. Proposition~\ref{thm:circulantnecessary} predicts this: below $16$, the only orders it allows are $1$ and $4$, and at both of these we already know examples.
\item For the odd lengths greater than $1$ there is also a simpler explanation, given in the pause at the end of Section~\ref{ch14:sec:periodic}: when $n$ is odd, every periodic correlation is odd, so none of them can be zero.
\item Every count in the table is even. By Exercise~\ref{ex:14.2}, negating every sign changes none of the correlations, aperiodic or periodic, because each correlation is a sum of pair products and pair products do not change (Section~\ref{ch06:sec:wlog}). So whenever a sequence $a$ qualifies, its negative $-a$ qualifies too. Moreover $-a$ is never equal to $a$, since that would require $a_0=-a_0$, that is, $a_0=0$. So the qualifying sequences of each length split into pairs $\{a,-a\}$, and their number is even.
\end{itemize}
The normalization of Section~\ref{ch06:sec:wlog} would allow a faster search: examine only the sequences that begin with $1$, and double each count. The program in Appendix~\ref{app:practical} examines all $2^n$ sequences instead, so the evenness of every count is a small independent check on its output.

\section{What current research adds}\label{ch14:sec:research}
The question of this chapter is a well-known research problem. The \emph{circulant Hadamard conjecture}, usually attributed to H.~J.~Ryser, states that there is no circulant Hadamard matrix of order greater than $4$. In the language of this chapter: apart from the lengths $1$ and $4$, no sign sequence has all of its periodic correlations at nonzero shifts equal to zero. The conjecture has been open since the 1960s. Proposition~\ref{thm:circulantnecessary} shows that a counterexample would need an order $4u^2$ with $u\ge2$. Over the decades, researchers have proved much stronger restrictions on the possible orders, but none of these results rules out all of them.

Barker sequences have a related history. Since the early 1960s it has been known that a Barker sequence of odd length has length at most $13$. In our table, the odd lengths that have Barker sequences are exactly $1$, $3$, $5$, $7$, $11$, and $13$. For even lengths, the Barker question turns into a question about circulant Hadamard matrices. It has long been known that when $n>2$ is even, every Barker sequence of length $n$ also has $C_s=0$ for every shift $s$ with $1\le s\le n-1$, so its circulant matrix is a Hadamard matrix. We do not prove this here, but you can see one instance of it in the table. The eight Barker sequences of length four are $(1,1,1,-1)$, its three other rotations, and the negatives of these four. They are also the eight sequences counted in the last column. As a consequence, a Barker sequence of even length greater than $4$ would give a circulant Hadamard matrix of order greater than $4$, and whether such a sequence exists was likewise an open question.

In September 2026, a preprint claiming a proof of the circulant Hadamard conjecture was released in OpenAI's public mathematics repository, together with a formalization in Lean~\cite{oai179}. Lean is a \emph{proof assistant}: a computer program that checks every step of a proof written in its own precise language. A \emph{formalization} of a theorem is a version of its statement and proof written in such a language, so that the program can check it. Chapters~\ref{ch:lean1} and~\ref{ch:lean2} introduce Lean.

According to the repository's description, the formalized statement says that a real circulant Hadamard matrix of positive order $n$ exists exactly when $n=1$ or $n=4$. The word ``real'' separates the matrices of this chapter, whose entries are the real numbers $1$ and $-1$, from a related family with complex entries that researchers also study. A second formalized statement, described there as a consequence, says that a sign sequence of positive even length whose aperiodic correlations at all nonzero shifts have absolute value at most one must have length $2$ or $4$. In our terms, every Barker sequence of even length has length $2$ or $4$. The description also notes that the preprint's treatment of odd lengths is not part of the formalized statements.

These claims come from an outside source, and nothing in this chapter establishes them. Our experiment stops at length $13$, and Proposition~\ref{thm:circulantnecessary} still permits the orders $16,36,64,\ldots$. What the chapter does give you is the ability to read the claims precisely. You know what a circulant matrix is, how the orthogonality of its rows translates into periodic correlations, why even-length Barker sequences belong to the same story, and exactly which orders an elementary argument leaves open. Ruling out the orders $4u^2$ with $u\ge2$ requires far more mathematics than this course covers. When you meet a summary of a result like this, whether in a news report or from an AI assistant, first check which statement was proved: which matrices, which kind of correlation, which lengths. Chapter~\ref{ch:bridge} returns to this example and describes how to begin reading a proof whose background you do not yet have.

\section*{Exercises}
\begin{exercise}\label{ex:14.1}For the sequence $(1,1,-1)$, compute the aperiodic correlations $A_1$ and $A_2$ and the periodic correlations $C_1$ and $C_2$ directly from their definitions. Check that your answers satisfy equation~\eqref{eq:periodaperiod}. Is $(1,1,-1)$ a Barker sequence?\end{exercise}
\begin{exercise}\label{ex:14.2}Let $a$ be a sign sequence of length $n$, and let $-a=(-a_0,-a_1,\ldots,-a_{n-1})$ be its negative. Prove that $A_s(-a)=A_s(a)$ for every shift $s$ with $1\le s\le n-1$, and that $C_s(-a)=C_s(a)$ for every shift $s$ with $0\le s\le n-1$.\end{exercise}
\begin{exercise}\label{ex:14.3}Let $H$ be the circulant matrix of $(1,1,1,-1)$ displayed in Section~\ref{ch14:sec:circulant}, with rows and columns numbered from $0$ to $3$. Compute the entries of $HH^{\mathsf T}$ in positions $(0,2)$, $(1,3)$, and $(2,2)$ directly, as dot products of rows. Check that they agree with $HH^{\mathsf T}=4I$, and for each of the two off-diagonal entries, name the periodic correlation $C_s$ that it equals.\end{exercise}
\begin{exercise}\label{ex:14.4}Let $a$ be a sign sequence of length $n$. Prove that $C_{n-s}=C_s$ by renaming the summation index in the sum that defines $C_{n-s}$. For which shifts $s$ are both sides of this equation defined?\end{exercise}
\begin{exercise}\label{ex:14.5}Use the two necessary conditions of this chapter (the order of a circulant Hadamard matrix is a perfect square, and a Hadamard matrix of order greater than $2$ has an order divisible by four) to decide which of the numbers $8$, $9$, $16$, $25$, $36$ could be the order of a circulant Hadamard matrix. Does your answer show that a circulant Hadamard matrix of any of these orders exists?\end{exercise}
\begin{exercise}\label{ex:14.6}A computer examines all $2^n$ sign sequences of one fixed length $n$, looking for sequences with $C_s=0$ for every shift $s$ with $1\le s\le n-1$. Explain why finding a single such sequence proves that a circulant Hadamard matrix of order $n$ exists, while finding none proves only that no circulant Hadamard matrix of this one order $n$ exists. What would it take to prove that there is no circulant Hadamard matrix of any order greater than $4$?\end{exercise}

